\documentclass[../../../include/open-logic-section]{subfiles}

\begin{document}

\olfileid{his}{set}{mythology}

\olsection{చరిత్ర కంటే పురాణమే ఎక్కువా?}

శతాబ్దానికి పైగా కాలం గడిచిన తరువాత ఈ సంఘటనలను వెనక్కి తిరిగి చూస్తున్నప్పుడు, ఆ తీర్పులను పరిశీలించకుండా నమ్మకూడదు. ఫలితాలు నిస్సందేహంగా కొత్తవీ, ఉత్సాహకరమైనవీ, ఆశ్చర్యకరమైనవీ. కానీ అవి నిజంగా ఎంత దిగ్భ్రాంతికరమైనవి? జ్యామితీయ అంతఃప్రజ్ఞపై ఆధారపడకూడదని అవి నిజంగానే చూపించాయా?

దిగ్భ్రాంతి విషయానికొస్తే, డెడెకిండ్‌కు కాంటర్ రాసిన ప్రసిద్ధ ``\emph{je le vois, mais je ne le crois
pas}'' అనే మాటను సందర్భం నుంచి బాగా వేరు చేసి తీసుకున్నారని \citet{Gouvea2011} ఎత్తిచూపారు. ఆ సందర్భంలో మరికొంత భాగం ఇది (\citeauthor{Gouvea2011} నుంచి ఉటంకింపు):
\begin{quote}
ఈ విషయంపై నా ఉత్సాహంతో మీ దయను, ఓర్పును ఇన్ని విధాలుగా కోరుతున్నందుకు క్షమించండి. ఇటీవల మీకు పంపిన విషయాలు నాకే అంత అనూహ్యంగా, అంత కొత్తగా అనిపిస్తున్నాయి; గౌరవనీయ మిత్రమా, వాటి సరైనదనంపై మీ నిర్ణయం పొందే వరకు నాకు మనశ్శాంతి లేదు. మీరు నాతో ఏకీభవించనంత వరకు నేను చెప్పగలిగేది ఒక్కటే: \emph{je le vois, mais je ne le crois pas.}
\end{quote}
తన ఫలితం ``అంత అనూహ్యం, అంత కొత్తది'' అని కాంటర్‌కు తెలుసు. కానీ ఆయన దాన్ని ఎప్పుడైనా \emph{నమ్మశక్యం కానిదిగా} భావించారనేది సందేహాస్పదం. \citeauthor{Gouvea2011} చెప్పినట్లు, తాను ఇచ్చిన నిరూపణను తనిఖీ చేయమని డెడెకిండ్‌ను కోరుతున్నారంతే.

జ్యామితీయ అంతఃప్రజ్ఞ విషయానికొస్తే: పియానో తన స్థలాన్ని నింపే వక్రరేఖను ఎలాంటి చిత్రాలూ లేకుండానే ప్రచురించారు. కానీ హిల్బర్ట్ తన వక్రరేఖను ప్రచురించినప్పుడు తన ఉద్దేశాన్ని వివరించారు: ``క్రింది జ్యామితీయ అంతఃప్రజ్ఞ సహాయాన్ని తీసుకుంటే'', పియానో ఫలితాన్ని అర్థం చేసుకోవడానికి పాఠకులకు స్పష్టమైన మార్గం ఇస్తానన్నారు; వెంటనే \olref[his][set][pathology]{sec}లో చూపిన వాటిలాంటి \emph{చిత్రాల} వరుసను జోడించారు.

మరింత సాధారణంగా: ప్రచురిత నిరూపణల్లో చిత్రాలు కొంత ఫ్యాషన్ తగ్గినా, విషయాలను నిరూపించేటప్పుడు గణితశాస్త్రజ్ఞులు \emph{తరచుగా} చిత్రాలను ఉపయోగిస్తారనే వాస్తవాన్ని విస్మరించలేం. (స్థూలంగా చెప్పాలంటే: జ్యామితీయ అంతఃప్రజ్ఞను ఎప్పుడు నమ్మవచ్చో మంచి గణితశాస్త్రజ్ఞులకు తెలుసు.)

సంక్షిప్తంగా: అతిశయ ప్రచారాన్ని నమ్మవద్దు; కనీసం పరిశీలించకుండా నమ్మవద్దు. దీని గురించి మరింత చదవడానికి \citet{Giaquinto2007} చూడండి.

\end{document}
